\chapter{Proposed Methodology and Architectural Paradigms}
\label{ch:methodology}

To transcribe full-page handwritten mathematics exam scripts, we designed, implemented, and compared two fundamentally different architectures:
\begin{enumerate}
    \item \textbf{Paradigm A (Modular Dual-Head Hybrid Pipeline):} A divide-and-conquer approach that separates the problem into clear, interpretable steps: image deskewing, Sauvola binarization, 1D vertical IoU line segmentation, an algorithmic matrix grid parser, a dual-head dispatcher (TrOCR for English text and our custom DenseNet-Coverage model for math), and a 2D directed acyclic graph (DAG) to sort lines into correct reading order.
    \item \textbf{Paradigm B (End-to-End Multimodal Vision-Language Model):} A holistic sequence-to-sequence model using Qwen2.5-VL-7B. We take advantage of its native dynamic resolution and 2D Rotary Position Embeddings (2D-RoPE), and we engineer an \emph{Assistant Prefill Locking} mechanism to stop the model from chatting and force it to transcribe raw ink verbatim.
\end{enumerate}

In this chapter, we explain the design, mathematical formulation, and implementation of both pipelines.

\section{Paradigm A: The Modular Dual-Head Hybrid Pipeline}
\label{sec:paradigm_a_overview}

Our modular pipeline processes a full exam page in five sequential stages, as shown in Figure \ref{fig:hybrid_pipeline_flow}:
\begin{equation}
\mathbf{I}_{\text{raw}} \xrightarrow{\text{Stage 1}} \mathbf{I}_{\text{deskew}} \xrightarrow{\text{Stage 2}} \{B_1, \dots, B_K\} \xrightarrow{\text{Stage 3}} \{\text{Type}_k\} \xrightarrow{\text{Stage 4}} \{\hat{\mathbf{y}}_k\} \xrightarrow{\text{Stage 5}} \mathbf{Y}_{\text{doc}}
\end{equation}

\begin{figure}[htbp]
\centering
\fbox{\parbox{0.92\textwidth}{
\vspace{0.3cm}
\centering
\textbf{[Paradigm A: Modular Dual-Head Hybrid Pipeline Architecture]}
\\[0.3cm]
\small
\begin{tabular}{ll}
\toprule
\textbf{Pipeline Stage} & \textbf{Algorithmic Operator / Function} \\
\midrule
\textbf{Input Tensor} & Raw Full-Page Scan $\mathbf{I}_{\text{raw}} \in \mathbb{R}^{H \times W \times C}$ \\
\textbf{Stage 1: Preprocessing} & Projection Deskewing ($\theta^*$) + Sauvola Adaptive Binarization ($T(x,y)$) \\
\textbf{Stage 2: Line Segmentation} & Horizontal Projections + 1D Vertical IoU NMS ($\theta_{\text{IoU}} = 0.35$) \\
\textbf{Stage 3: Delimiter Extraction} & \texttt{MatrixGridParser} Aspect Filtering ($h/w > 3.0$) $\to$ \texttt{bmatrix} \\
\textbf{Stage 4: Dual Dispatcher} & \textbf{Text Head:} TrOCR ViT \quad \textbf{Math Head:} DenseNet-Coverage HMER \\
\textbf{Stage 5: Document Assembly} & 2D Spatial Reading-Order DAG $\mathcal{G} = (\mathcal{V}, \mathcal{E}, \prec)$ \\
\midrule
\textbf{Output Document} & Standalone Compilable \LaTeX\ Document ($\text{Compile}(\mathbf{Y}) = 1$) \\
\bottomrule
\end{tabular}
\vspace{0.3cm}
}}
\caption{Dataflow diagram showing the modular stages in our Hybrid Pipeline.}
\label{fig:hybrid_pipeline_flow}
\end{figure}

\section{Image Preprocessing and Adaptive Local Binarization}
\label{sec:preprocessing}

Raw scanner images often have minor rotational tilt from the scanner feed and uneven lighting caused by paper crinkles or shadows near page folds. Our preprocessing cleans up these issues while keeping delicate strokes intact.

\subsection{Deskewing via Projection Profile Variance}
To straighten tilted pages without using brittle Hough transforms on broken handwriting, we maximize the variance of horizontal projection profiles.

We test a range of rotation angles $\theta \in [-15^\circ, +15^\circ]$. For each angle, we rotate the image and calculate its horizontal projection histogram $H_\theta(y)$. When the angle $\theta^*$ matches the actual tilt of the handwriting, the text lines line up horizontally, creating tall peaks separated by empty zero-intensity valleys. This maximizes the variance:
\begin{equation}
\theta^* = \arg\max_{\theta \in [-15^\circ, +15^\circ]} \text{Var}\left( H_\theta(y) \right)
\end{equation}
We then rotate the page by $-\theta^*$ using bicubic interpolation to straighten all lines.

\subsection{Adaptive Local Sauvola Binarization}
Global binarization methods like Otsu's thresholding \cite{otsu1979threshold} pick one threshold for the whole page. This fails when students use different pens, or when lighting falls off toward the edges. To handle varying pen pressure and faint pencil marks, we use \textbf{Sauvola Adaptive Binarization} \cite{sauvola2000adaptive}. For each pixel $(x, y)$, we compute a local threshold $T(x, y)$ over a $31 \times 31$ pixel window:
\begin{equation}
T(x, y) = m(x, y) \cdot \left[ 1 + k \cdot \left( \frac{s(x, y)}{R} - 1 \right) \right]
\end{equation}
where $m(x, y)$ and $s(x, y)$ are the local mean and standard deviation of pixel intensities, $R = 128$ is the dynamic range for 8-bit grayscale, and $k = 0.20$ controls sensitivity to faint strokes. Pixels darker than $T(x, y)$ are marked as handwriting.

\section{Line Segmentation via 1D Vertical IoU}
\label{sec:iou_nms_segmentation}

Once binarized, the page needs to be cut into separate lines of text, formulas, and matrices without cutting through letters.

\subsection{Horizontal Projection Histograms}
We compute the horizontal projection profile along the vertical axis:
\begin{equation}
H(y) = \sum_{x=1}^{W} \mathbf{I}_{\text{bin}}(x, y), \quad y \in \{1, \dots, H\}
\end{equation}
We find candidate lines by looking for vertical bands where the projection density rises above an adaptive noise threshold ($H(y) > 0.02 \cdot \max_y H(y)$).

\subsection{1D Vertical IoU Clustering}
On unruled exam paper, students frequently write lines that interlock. The descender of a `$y$' on line 1 overlaps vertically with the ascender of a `$d$' on line 2. A simple projection cut would slice right through both letters.

We solve this by framing line grouping as \textbf{1D Vertical IoU Non-Maximum Suppression}. For each detected candidate slice $k$, we define its vertical interval as $J_k = [y_{k, \text{top}}, y_{k, \text{bottom}}]$:
\begin{enumerate}
    \item For any two overlapping candidate intervals $J_a$ and $J_b$, we calculate their 1D Vertical Intersection-over-Union:
    \begin{equation}
    \text{IoU}_{1D}(J_a, J_b) = \frac{\max\left(0, \ \min(y_{a, \text{bottom}}, y_{b, \text{bottom}}) - \max(y_{a, \text{top}}, y_{b, \text{top}})\right)}{\max(y_{a, \text{bottom}}, y_{b, \text{bottom}}) - \min(y_{a, \text{top}}, y_{b, \text{top}})}
    \end{equation}
    \item If $\text{IoU}_{1D}(J_a, J_b) \ge \theta_{\text{IoU}} = 0.35$, the two intervals belong to the same interlocking handwriting line. We merge them into a single bounding box.
    \item If $\text{IoU}_{1D}(J_a, J_b) < \theta_{\text{IoU}}$, we keep them as separate lines.
\end{enumerate}
This ensures tall characters remain intact inside their respective line crops.

\section{The Algorithmic Matrix Parser (\texttt{MatrixGridParser})}
\label{sec:matrix_grid_parser}

A major headache with standard OCR tools on linear algebra scripts is that they slice large matrices into separate, disconnected rows of numbers. A $4 \times 4$ matrix takes up four vertical lines of text enclosed by a tall pair of brackets. If a line segmenter cuts through those brackets, the matrix structure is destroyed.

To preserve matrices, we wrote an algorithmic parser called \texttt{MatrixGridParser}:
\begin{enumerate}
    \item \textbf{Bracket Detection:} We run connected-component analysis to find tall, slender vertical strokes that have an aspect ratio $h/w > 3.0$ and span more than two standard line heights. If we find matching left and right brackets bounding a region, we mark that bounding box as a matrix.
    \item \textbf{Row and Column Splitting:} Inside the matrix box, we compute vertical and horizontal projection profiles to identify the rows and columns of individual cell entries.
    \item \textbf{Cell Recognition and Assembly:} Each isolated cell crop is recognized, and we wrap the entire grid inside a \LaTeX\ \texttt{\textbackslash begin\{bmatrix\} \dots \textbackslash end\{bmatrix\}} block, using \texttt{\&} to separate columns and \texttt{\textbackslash\textbackslash} to separate rows.
\end{enumerate}
This completely bypasses the risk of line segmenters cutting through matrix brackets.

\section{Dual-Head Recognition Dispatcher}
\label{sec:dual_head_dispatcher}

Once the page is segmented into bounding boxes, each box must be transcribed. Handwritten exam scripts contain both cursive English proofs (*"Assume $W$ is a subspace..."*) and formal mathematical formulas. Using a text OCR model on math produces garbage, and using a math formula model on English sentences causes it to crash.

We built a classifier to route each crop to the right model:
\begin{itemize}
    \item \textbf{Text Head (TrOCR):} Crops classified as English prose are sent to Microsoft's \texttt{trocr-base-handwritten} Vision Transformer \cite{li2023trocr}.
    \item \textbf{Math Head (DenseNet-Coverage HMER):} Crops containing mathematical formulas are routed to our custom 23.1-million parameter DenseNet-Coverage model (\texttt{src/models/hmer\_model.py}), which uses an Attention Refinement Module to track formula syntax.
\end{itemize}
Our classifier examines stroke density, aspect ratio, and connected components to make the routing decision.

\section{Reading-Order Recovery via 2D Spatial DAG}
\label{sec:spatial_dag}

After all boxes are transcribed, they must be arranged in the correct reading order. Students often write margin notes or place equations side-by-side. 

We represent the page layout as a Directed Acyclic Graph (DAG) $\mathcal{G} = (\mathcal{V}, \mathcal{E}, \prec)$, where each vertex $v \in \mathcal{V}$ is a recognized bounding box. A directed edge $(u, v) \in \mathcal{E}$ indicates that block $u$ must be read before block $v$. The ordering relation $u \prec v$ is determined by:
\begin{equation}
u \prec v \iff \begin{cases}
y_{u, \text{bottom}} < y_{v, \text{top}} & \text{(Block $u$ is physically above block $v$)} \\
x_{u, \text{right}} < x_{v, \text{left}} \quad \text{when } \text{Overlap}_y(u, v) > 0.5 & \text{(Block $u$ is to the left of block $v$)}
\end{cases}
\end{equation}
A standard topological sort on this DAG gives the final document reading order, which we write to \texttt{output\_solution.tex}.

\section{Paradigm B: The Multimodal Vision-Language Frontier}
\label{sec:paradigm_b_overview}

While our Modular Hybrid Pipeline gives full control over every step, errors in early stages (like line segmentation) can ruin later stages. This led us to explore an alternative: an end-to-end vision-language model, specifically **Qwen2.5-VL-7B-Instruct** \cite{wang2024qwen2vl}.

Qwen2.5-VL takes the entire unsegmented page image as input and generates a complete \LaTeX\ document in a single autoregressive pass. It uses dynamic visual resolution and 2D Rotary Position Embeddings (2D-RoPE) \cite{su2024roformer} to understand 2D layouts directly in self-attention.

\subsection{The Conversational Collapse Problem}
However, when we first ran Qwen2.5-VL zero-shot on exam scans, it failed in a frustrating way. Because the model was aligned with RLHF to be a helpful conversational assistant, it refused to just output \LaTeX\ code. Instead, it started chatting:
\begin{quote}
\emph{``Certainly! Below is the transcription of the handwritten linear algebra exam script shown in the image...''}
\end{quote}
Once the model started generating conversational English, it stopped paying attention to the image and failed to produce a compilable document.

\subsection{Our Solution: Assistant Prefill Locking}
To stop this conversational drift permanently, we designed **Assistant Prefill Locking** using Hugging Face's \texttt{continue\_final\_message=True} flag.

In Assistant Prefill Locking, we prepopulate the assistant's response buffer with a formal \LaTeX\ document preamble before generation begins:
\begin{equation}
\label{eq:prefill_string}
\mathbf{Y}_{\text{prefill}} = \big[ \texttt{\textbackslash documentclass[12pt]\{article\}}, \dots, \texttt{\textbackslash begin\{document\}}, \dots, \texttt{\textbackslash section*\{Solution\}} \big]
\end{equation}

The full prompt fed to the vision-language model is structured as:
\begin{equation}
\label{eq:prompt_structure}
\begin{aligned}
\mathcal{T}_{\text{input}} = \big[ &\texttt{<|im\_start|>system} \ \dots \ \texttt{<|im\_end|>}, \\
&\texttt{<|im\_start|>user} \ [\mathbf{X}_{\text{vis}}; \ \text{Prompt}] \ \texttt{<|im\_end|>}, \\
&\texttt{<|im\_start|>assistant} \ \mathbf{Y}_{\text{prefill}} \big]
\end{aligned}
\end{equation}

When inference starts with \texttt{continue\_final\_message=True}, the model cannot generate an assistant introduction or say "Certainly!". Its very first output token must continue immediately from the end of $\mathbf{Y}_{\text{prefill}}$ at token position $T_{\text{prefill}} + 1$:
\begin{equation}
y_{T_{\text{prefill}} + 1} \sim P\left( \cdot \mid \mathbf{Y}_{\text{prefill}}, \ \mathbf{X}_{\text{vis}} \right)
\end{equation}
Because the context buffer ends right after \texttt{\textbackslash section*\{Solution\}\textbackslash n\textbackslash n}, the chance of generating conversational filler drops to zero. The model is forced to output mathematical \LaTeX\ tokens immediately, diving straight into transcribing the visual ink.

\section{Negative Constraints and Regex Cleaners}
\label{sec:negative_constraints}

To make sure the model transcribes literal handwriting rather than hallucinating solutions, we use a two-part safeguard:
\begin{enumerate}
    \item \textbf{Negative System Prompt:} In \texttt{QwenVLOCREngine}, our system prompt forbids the model from solving the math, making corrections, or outputting Python scripts (\texttt{import numpy as np}). We explicitly tell the model: *"You are an OCR transcription engine, not a math tutor. Transcribe exactly what is written, mistakes included."*
    \item \textbf{Regex Post-Processing:} We run regex filters on the generated text to strip away any residual Markdown code fences (\texttt{```latex}) and remove leaked Python snippets.
\end{enumerate}

\section{Repetition Penalty Calibration and the Math Paradox}
\label{sec:repetition_penalty}

In linear algebra, students frequently perform row operations that produce consecutive zeros or repeated equation patterns. In greedy decoding ($\text{temperature} = 0.0$), the attention mechanism can fall into a loop, repeatedly outputting zeros until it hits the token limit.

\subsection{The NLP Repetition Penalty Paradox}
Standard NLP practice prevents loops by applying a strong repetition penalty ($\rho > 1.2$) or blocking repeated n-grams (\texttt{no\_repeat\_ngram\_size = 15}). 

However, during our tests on PARAM Kamrupa, we discovered the **Repetition Penalty Paradox in Mathematical Transcription**. In math, repetition across consecutive lines is normal and necessary:
\begin{itemize}
    \item \textbf{Repeated Coordinate Equations:} When checking if a matrix is in a subspace, a student writes out four similar coordinate equations: $\sum_{j=1}^4 a_{ij} = 0$ for $i \in \{1, 2, 3, 4\}$.
    \item \textbf{Consecutive Row Operations:} Students perform multiple row operations that start with the same notation (e.g., $R_2 \to R_2 - 2R_1$, followed by $R_3 \to R_3 - 2R_1$).
\end{itemize}
When we had \texttt{no\_repeat\_ngram\_size = 15} enabled, the beam search artificially blocked valid subsequent equations because their syntax looked identical to the previous line. This was the exact reason why four consecutive coordinate equations were skipped on Page 1 of Roll 210103096 in our early runs.

\subsection{Fixing Loops with Matrix Regex Collapsing}
To fix repetitive loops without blocking authentic math, we changed our approach:
\begin{enumerate}
    \item \textbf{Low Global Penalty:} We set a modest repetition penalty ($\rho = 1.05$--$1.10$) and disabled multi-line n-gram blocking completely.
    \item \textbf{Regex Matrix Collapser:} If the model gets stuck generating repeating zeros inside a matrix environment (\texttt{\textbackslash begin\{bmatrix\} \dots \textbackslash end\{bmatrix\}}), our post-processor detects when consecutive zeros exceed the column rank ($N > 6$) and trims the row back to the correct matrix dimension.
\end{enumerate}

\section{The Pushdown Stack Balancer for 100\% Compilation}
\label{sec:environment_balancer}

A strict requirement of this project is that every single generated \LaTeX\ file must compile into a PDF without crashing: $\text{Compile}(\hat{\mathbf{Y}}) = 1$.

In autoregressive generation, if a student writes a very long answer, the model might hit its maximum token limit ($T_{\text{max}} = 2560$) midway through an equation or matrix. If generation cuts off inside an open \texttt{\textbackslash begin\{bmatrix\}}, running \texttt{pdflatex} causes a fatal syntax crash.

To prevent syntax crashes under all conditions, we built a deterministic **Pushdown Automaton (Stack Balancer)** inside \texttt{src/latex/generator.py}:
\begin{enumerate}
    \item We initialize an empty stack $S$. Whenever the parser encounters \texttt{\textbackslash begin\{env\_name\}}, it calls $\text{Push}(S, \text{env\_name})$.
    \item Whenever it encounters \texttt{\textbackslash end\{env\_name\}}, if $\text{Top}(S) == \text{env\_name}$, it calls $\text{Pop}(S)$.
    \item When generation terminates, if $S$ is not empty, we pop each open environment and append its closing tag:
    \begin{equation}
    \hat{\mathbf{Y}} \leftarrow \hat{\mathbf{Y}} \ \oplus \ \texttt{\textbackslash end\{} e \texttt{\}\textbackslash n}
    \end{equation}
    \item Finally, if \texttt{\textbackslash end\{document\}} is missing, we append it:
    \begin{equation}
    \hat{\mathbf{Y}} \leftarrow \hat{\mathbf{Y}} \ \oplus \ \texttt{\textbackslash end\{document\}}
    \end{equation}
\end{enumerate}
Together with our prefilled preamble, this stack balancer mathematically guarantees that every output file is syntactically closed, giving us a **$100\%$ Compilation Pass Rate**.

\section{Comparing the Two Paradigms}
\label{sec:methodological_synthesis}

Table \ref{tab:paradigm_comparison} compares our Modular Hybrid Pipeline with the Multimodal VLM across six key practical dimensions.

\begin{table}[htbp]
\centering
\caption{Comparison Between Our Modular Hybrid Pipeline and Multimodal VLM}
\label{tab:paradigm_comparison}
\vspace{0.2cm}
\small
\begin{tabular}{p{2.8cm}p{5.2cm}p{5.2cm}}
\toprule
\textbf{Dimension} & \textbf{Paradigm A: Modular Hybrid Pipeline} & \textbf{Paradigm B: Multimodal VLM (Qwen2.5-VL)} \\
\midrule
\textbf{Layout Dependency} & High: Needs explicit 1D IoU line segmentation and contour detection. & None: Full-page visual cross-attention parses layout directly. \\
\addlinespace
\textbf{2D Matrix Parsing} & Explicit algorithmic projection via \texttt{MatrixGridParser} into \texttt{bmatrix}. & Native 2D spatial attention using 2D-RoPE embeddings. \\
\addlinespace
\textbf{Text vs. Math Routing} & Heuristic classifier routes crops to TrOCR (Text) vs. DenseNet (Math). & Single unified model processes cursive text and formulas together. \\
\addlinespace
\textbf{Hardware Requirements} & Lightweight ($\approx 23.1\text{M}$ params HMER + $334\text{M}$ TrOCR); runs easily on CPU/GPU. & Heavy ($\approx 7.6\text{B}$ params); needs a dedicated NVIDIA Tesla V100 GPU (16GB VRAM). \\
\addlinespace
\textbf{Hallucination Risk} & Zero language hallucination; strictly bound to cropped pixel boxes. & Moderate without prompts; eliminated using Assistant Prefill Locking. \\
\addlinespace
\textbf{Compilation Invariant} & Combines independently recognized strings into a template. & Stack-based pushdown automaton closes truncated environments. \\
\bottomrule
\end{tabular}
\end{table}

\section{Paradigm C: The Master Solution --- Multi-Band High-Resolution Slicing Architecture}
\label{sec:multiband_slicing}

While Paradigm B demonstrated the representational power of unified vision-language modeling, deploying a 7-billion parameter multimodal model on uncompressed, full-page exam scripts encountered a fundamental hardware barrier on the PARAM Kamrupa Supercomputer. 

\subsection{The Optical Resolution vs. Quadratic ViT Memory Dilemma}
In a Vision Transformer (ViT), the 2D optical image $I \in \mathbb{R}^{H \times W \times C}$ is partitioned into non-overlapping spatial patches of dimension $P \times P$ (where $P = 14\,\text{px}$). The total number of visual tokens fed to the transformer encoder is given by:
\begin{equation}
N_{\text{patches}} = \left\lfloor \frac{H}{P} \right\rfloor \times \left\lfloor \frac{W}{P} \right\rfloor
\end{equation}
An uncompressed 300 DPI examination sheet ($1653 \times 2338\,\text{px}$) yields:
\begin{equation}
N_{\text{patches}} \approx \left\lfloor \frac{2338}{14} \right\rfloor \times \left\lfloor \frac{1653}{14} \right\rfloor \approx 167 \times 118 \approx 19,706 \text{ raw patches}
\end{equation}
Even under native sub-sampling, an unconstrained full-page input generates over $5,800$ image patches. The multi-head self-attention mechanism computes the full pairwise attention score matrix:
\begin{equation}
\text{Attention}(Q, K, V) = \text{Softmax}\left(\frac{Q K^T}{\sqrt{d_k}}\right) V
\end{equation}
The attention tensor $A = Q K^T$ requires memory scaling quadratically with the spatial token sequence length:
\begin{equation}
\text{Memory}_{\text{attention}} = B \times H_{\text{heads}} \times N_{\text{patches}}^2 \times 2\,\text{bytes}
\end{equation}
For $N_{\text{patches}} \approx 5,800$, allocating this contiguous memory block requires approximately $1.30\,\text{GiB}$ in a single contiguous allocation. On an NVIDIA Tesla V100 GPU with a 16 GB VRAM ceiling, the model weights in FP16 precision consume $14.07\,\text{GiB}$, leaving only $1.71\,\text{GiB}$ of addressable headroom. Because CUDA memory becomes fragmented by runtime buffers, the memory allocator fails to locate a contiguous $1.30\,\text{GiB}$ block, throwing a fatal \texttt{CUDA Out of Memory} crash.

Conversely, downscaling the full page to fit within the GPU memory headroom forces the image width down to $848\,\text{px}$. Compressing 30 handwritten lines into $848\,\text{px}$ results in catastrophic optical degradation: fine pen strokes like minus signs ($-$), decimal points ($.$), subscripts ($a_i$), and primes ($E'$) blur into the paper background. A critical coefficient like $3.5$ is frequently misrecognized as $8.5$, rendering mathematical transcription unviable.

\subsection{The 3-Band Overlapping Partitioning Algorithm}
To solve both the optical resolution dilemma and the GPU memory ceiling simultaneously, we engineered the **Multi-Band High-Resolution Slicing Architecture**. Human evaluators do not grade an entire 30-line answer booklet in a single glance; they read it section by section, derivation by derivation. We modeled this visual attention mechanism computationally through an automated 3-stage pipeline:

\begin{enumerate}
    \item \textbf{Stage 1: Booklet Header Instruction Stripping:} The top $18\%$ of every IIT Guwahati answer booklet contains printed static instructions, barcodes, and candidate roll-number boxes ($y \le 0.18H$). The pipeline automatically detects and crops this static banner:
    \begin{equation}
    I_{\text{active}} = I\left[0.18H : H,\, 0 : W\right]
    \end{equation}
    Eliminating the printed header prevents the vision-language model from wasting prompt tokens on irrelevant institutional instructions.
    
    \item \textbf{Stage 2: 3-Band Overlapping Partition with Safety Buffer:} The remaining active answer canvas ($H_{\text{active}} \approx 0.82H$) is divided into three horizontal bands ($B_1, B_2, B_3$). Crucially, to prevent boundary expressions from being fractured, we enforce an $8\%$ vertical safety overlap buffer ($\delta \approx 180\,\text{px}$):
    \begin{align}
    B_1 &= I_{\text{active}}\left[0.00H_{\text{active}} : 0.38H_{\text{active}} + \delta,\, 0 : W\right] \\
    B_2 &= I_{\text{active}}\left[0.34H_{\text{active}} - \delta : 0.72H_{\text{active}} + \delta,\, 0 : W\right] \\
    B_3 &= I_{\text{active}}\left[0.68H_{\text{active}} - \delta : 1.00H_{\text{active}},\, 0 : W\right]
    \end{align}
    Because handwritten derivation lines in MA102 booklets average $80$--$120\,\text{px}$ in height, an overlap buffer of $\delta \approx 180\,\text{px}$ mathematically guarantees that any mathematical fraction, matrix row, or integral sign crossing a partition cut appears completely intact in at least one adjacent slice.
    
    \item \textbf{Stage 3: High-Resolution Normalization:} Each band is processed independently at a native dimension of $\text{max\_dim} = 1200\,\text{px}$ (yielding an optical raster of $\approx 1200 \times 470\,\text{px}$). Because each band contains only $7$--$8$ lines of student handwriting instead of 30, the effective token density per line increases by $3.5\times$ compared to monolithic processing.
\end{enumerate}

\subsection{Quadratic Attention Memory Collapse}
By partitioning the canvas into 3 horizontal slices, the visual token sequence length per forward pass is reduced from $N \approx 5,800$ to $N \approx 1,900$ patches. Since transformer self-attention scales quadratically ($N^2$), the peak activation memory collapses:
\begin{equation}
\text{Memory}_{\text{slice}} \approx \left(\frac{1}{3}\right)^2 \times \text{Memory}_{\text{monolithic}} \approx \frac{1}{9} \times 1.30\,\text{GiB} \approx 240\,\text{MiB}
\end{equation}
This drops the peak GPU memory to under $14.35\,\text{GiB}$, leaving over $1.43\,\text{GiB}$ of guaranteed safe headroom on our 16 GB Tesla V100 GPU, permanently eliminating CUDA Out-of-Memory crashes.

\section{Automated LaTeX Stitching and Direct Compilation Pipeline}
\label{sec:latex_stitching_pipeline}

Once the three horizontal image slices ($B_1, B_2, B_3$) are independently transcribed by the multimodal model, we face a critical systems engineering problem: how do we stitch three separate, asynchronous text fragments into a single, syntactically unified \LaTeX\ document that compiles with zero human intervention?

Our post-processing engine (\texttt{stitch\_latex\_parts.py}) executes four automated algorithmic phases:

\begin{figure}[htbp]
\centering
\begin{minipage}{0.95\textwidth}
\begin{center}
\fbox{
\begin{minipage}{0.92\textwidth}
\vspace{0.2cm}
\textbf{Algorithm: Automated Four-Stage LaTeX Stitching and Compilation Engine}
\begin{enumerate}[leftmargin=1.5em]
    \item \textbf{Body Sanitization \& Prefill Extraction:} Extract document body between \texttt{\textbackslash begin\{document\}} and \texttt{\textbackslash end\{document\}}, stripping duplicate preambles and markdown code fences (\texttt{```latex}).
    \item \textbf{Suffix-Prefix Boundary Deduplication:} Compute character-level normalized Levenshtein similarity across the $8\%$ overlap boundary to eliminate repeated derivation lines.
    \item \textbf{Pushdown Automaton (DPDA) Environment Balancer:} Maintain a LIFO stack $S$ tracking open \texttt{\textbackslash begin\{env\}} tags across slice cuts; synthesize matching \texttt{\textbackslash end\{env\}} tags in reverse topological order.
    \item \textbf{Underscore Regex Masking Engine:} Mask all valid math blocks via hash tokens \texttt{ZZMATHBLOCKiZZ}, escape prose shorthand underscores (\texttt{\textbackslash\_}), and demask math.
    \item \textbf{Autonomous Tectonic XeTeX Compilation:} Compile stitched code to PDF via embedded non-interactive Tectonic compiler in $< 5$ seconds.
\end{enumerate}
\vspace{0.1cm}
\end{minipage}
}
\end{center}
\end{minipage}
\caption{Algorithmic Flow of the Automated LaTeX Stitching Pipeline}
\label{fig:stitching_pipeline}
\end{figure}

\subsection{Formal Grammar Formulation: The Pushdown Automaton Balancer}
In formal language theory, \LaTeX\ mathematical markup forms a \textbf{Deterministic Context-Free Language (DCFL)} generated by a context-free grammar $G = (V, \Sigma, R, S)$ \cite{hopcroft2001introduction}. Delimited environments like \texttt{\textbackslash begin\{bmatrix\}} and \texttt{\textbackslash end\{bmatrix\}}, or \texttt{\textbackslash begin\{aligned\}} and \texttt{\textbackslash end\{aligned\}}, define nested Dyck scopes of arbitrary depth that cannot be parsed by a regular finite automaton.

When an answer script is sliced into horizontal bands, a partition boundary cut can occur midway through an open mathematical environment. If the generated slice terminates with an unclosed \texttt{\textbackslash begin\{bmatrix\}}, the entire document will fail to compile. To guarantee syntax closure under all slicing conditions, we formulated a \textbf{Deterministic Pushdown Automaton (DPDA)}:
\begin{equation}
M = \left( Q, \Sigma, \Gamma, \delta, q_0, Z_0, F \right)
\end{equation}
where:
\begin{itemize}
    \item $Q = \{ q_{\text{scan}}, q_{\text{open}}, q_{\text{close}}, q_{\text{repair}}, q_{\text{accept}} \}$ is the finite set of internal states.
    \item $\Sigma$ is the input alphabet consisting of ASCII and UTF-8 \LaTeX\ tokens, control sequences, and environment delimiters.
    \item $\Gamma = \{ Z_0, \texttt{bmatrix}, \texttt{pmatrix}, \texttt{aligned}, \texttt{align*}, \texttt{equation*}, \texttt{cases}, \dots \}$ is the stack alphabet, with $Z_0$ denoting the bottom-of-stack marker.
    \item $q_0 = q_{\text{scan}}$ is the initial start state with initial stack configuration $[Z_0]$.
    \item $F = \{ q_{\text{accept}} \}$ is the set of accepting terminal states.
\end{itemize}

The deterministic transition function $\delta: Q \times (\Sigma \cup \{\varepsilon\}) \times \Gamma \to Q \times \Gamma^*$ executes according to formal transition rules:
\begin{enumerate}
    \item \textbf{Environment Push ($\delta_{\text{push}}$):} Upon scanning an opening tag $\texttt{\textbackslash begin\{}e\texttt{\}}$ for any $e \in \Gamma$, the automaton pushes $e$ onto the LIFO stack:
    \begin{equation}
    \delta(q_{\text{scan}}, \texttt{\textbackslash begin\{}e\texttt{\}}, X) = (q_{\text{scan}}, eX), \quad \forall X \in \Gamma
    \end{equation}
    \item \textbf{Environment Pop ($\delta_{\text{pop}}$):} Upon scanning a closing tag $\texttt{\textbackslash end\{}e\texttt{\}}$, if the top of the stack matches $e$, the automaton pops $e$:
    \begin{equation}
    \delta(q_{\text{scan}}, \texttt{\textbackslash end\{}e\texttt{\}}, e) = (q_{\text{scan}}, \varepsilon)
    \end{equation}
    \item \textbf{EOF Stack Unwind ($\delta_{\text{repair}}$):} Upon reaching end-of-file ($\text{EOF}$) while the stack contains unpopped environments ($\text{Top}(S) = e \neq Z_0$), the automaton executes $\varepsilon$-transitions, synthesizing matching closing delimiters in reverse topological order:
    \begin{equation}
    \delta(q_{\text{repair}}, \varepsilon, e) = (q_{\text{repair}}, \varepsilon) \quad \text{emitting } \texttt{\textbackslash end\{}e\texttt{\}}
    \end{equation}
    \item \textbf{Acceptance:} When $\text{Top}(S) = Z_0$ at $\text{EOF}$, the automaton transitions to $q_{\text{accept}}$ and emits $\texttt{\textbackslash end\{document\}}$.
\end{enumerate}
This deterministic automaton guarantees that regardless of where physical cuts occur, the synthesized Abstract Syntax Tree is mathematically balanced and compiles with a $100\%$ pass rate.

\subsection{The Underscore Masking Regex Engine}
In \LaTeX, the underscore character \texttt{\_} is a reserved mathematical subscript operator. When an unescaped underscore appears in regular English text, the \LaTeX\ compiler halts with a fatal \texttt{Missing \$ inserted} syntax error. In student exam papers, students frequently use shorthand variable notations in English sentences (such as \texttt{fourth\_col} or \texttt{R\_1\_new}).

Applying a naive global regex replace (\texttt{s/\_/\textbackslash\_/g}) corrupts all valid mathematical subscript expressions (turning $W_1$ into $W\backslash\_1$). To solve this, our engine implements a three-step masking mechanism:
\begin{enumerate}
    \item \textbf{Math Block Extraction:} All inline math (\texttt{\$...\$}) and display math (\texttt{\textbackslash[...\textbackslash]}, \texttt{\textbackslash begin\{equation\}...\textbackslash end\{equation\}}) blocks are extracted and replaced with unique alphanumeric hash tokens: \texttt{ZZMATHBLOCK0ZZ}, \texttt{ZZMATHBLOCK1ZZ}, $\dots$
    \item \textbf{Prose Underscore Escaping:} A global regex scans the remaining document prose, escaping all bare underscores: \texttt{\_} $\to$ \texttt{\textbackslash\_}.
    \item \textbf{Demasking Restoration:} The unique hash tokens are iteratively substituted back with their original mathematical blocks untouched.
\end{enumerate}

\subsection{Autonomous Tectonic XeTeX Compilation}
Finally, the unified document is written to disk and passed directly to the embedded **Tectonic (XeTeX)** compiler:
\begin{verbatim}
tectonic -X compile output_stitched.tex
\end{verbatim}
Tectonic is a modern, standalone TeX engine that runs non-interactively, automatically caches font packages, and completes multi-pass references in memory. The entire compilation cycle from stitched \LaTeX\ source to publication-grade PDF executes in **under 5 seconds with a 100\% compilation pass rate**.

\section{Comprehensive Architectural Synthesis: Monolithic vs. Multi-Band Slicing}
\label{sec:architectural_synthesis_matrix}

Table \ref{tab:monolithic_vs_slicing_comparison} presents a comprehensive architectural comparison across seven critical engineering dimensions, demonstrating why Multi-Band High-Resolution Slicing represents the optimal paradigm for full-page exam script transcription.

\begin{table}[htbp]
\centering
\caption{Comprehensive Comparison: Monolithic Full-Page vs. Multi-Band Slicing}
\label{tab:monolithic_vs_slicing_comparison}
\vspace{0.2cm}
\small
\begin{tabular}{p{3.2cm}p{3.8cm}p{3.8cm}p{2.4cm}}
\toprule
\textbf{Architectural Metric} & \textbf{Monolithic Full-Page} & \textbf{Multi-Band Slicing} & \textbf{Engineering Gain} \\
\midrule
\textbf{ViT Input Patches ($N$)} & $> 5,800$ patches / page & $\approx 1,900$ patches / slice & $67\%$ reduction \\
\addlinespace
\textbf{ViT Attention Memory} & $1.30\,\text{GiB}$ (contiguous) & $\approx 240\,\text{MiB}$ & $81.5\%$ VRAM drop \\
\addlinespace
\textbf{GPU Safety Headroom} & $< 450\,\text{MiB}$ (High OOM risk) & $> 1,400\,\text{MiB}$ (Safe buffer) & Zero CUDA OOM \\
\addlinespace
\textbf{Effective Pixel Width} & $848\,\text{px}$ (downscaled) & $1200\,\text{px}$ (native slice) & $1.41\times$ wider raster \\
\addlinespace
\textbf{Lines per Visual Frame} & $25$--$30$ handwritten lines & $7$--$8$ handwritten lines & $3.5\times$ token density \\
\addlinespace
\textbf{Subscripts / Minus Signs} & Blurry; $-$, $.$ dropped & Crisp stroke definition & High optical clarity \\
\addlinespace
\textbf{Token-Level Accuracy} & $\approx 60\%$ (dropped proofs) & \textbf{88.24\%} (Levenshtein TER) & \textbf{+28.2\% accuracy} \\
\bottomrule
\end{tabular}
\end{table}
